\documentclass[11pt]{article}
\usepackage{amsmath,amssymb,amsthm}
\usepackage[margin=1in]{geometry}
\usepackage[colorlinks=true,linkcolor=blue]{hyperref}

\newtheorem{conjecture}{Conjecture}
\newtheorem{claim}{Claim}
\theoremstyle{definition}
\newtheorem{remark}{Remark}

\title{Disproof of TLMC Conjecture 00000000325\\
\large The packing dimension of $W(\tau)$ is not $2/(1+\tau)$}
\author{}
\date{}

\begin{document}
\maketitle

\section{The conjecture}

\begin{conjecture}[TLMC 00000000325]
Let
\[
W(\tau) \;=\; \bigl\{\alpha : \liminf_{q\to\infty}\; q^{1/\tau}\,\|q\alpha\| = 0\bigr\}.
\]
Then the packing dimension of $W(\tau)$ is exactly $2/(1+\tau)$.
\end{conjecture}

\section{Verdict: FALSE}

The conjecture is false. A single explicit counterexample suffices.

\begin{claim}[Primary counterexample: $\tau = \tfrac12$]
At $\tau = \tfrac12$ the value asserted by Conjecture~1 is
\[
\frac{2}{1+\tau} \;=\; \frac{2}{1+\tfrac12} \;=\; \frac{4}{3} \;>\; 1,
\]
which exceeds the dimension of the ambient space $\mathbb{R}$.
\end{claim}

\begin{proof}
Arithmetic: $2/(1+1/2) = 2/(3/2) = 4/3$, and $4 \cdot 1 = 4 > 3 = 1\cdot 3$,
so $4/3 > 1$. But the packing dimension is monotone under inclusion and
$\dim_P(\mathbb{R}) = 1$; hence $\dim_P(A) \le 1$ for \emph{every}
$A \subseteq \mathbb{R}$. No set can have packing dimension $4/3$, so the
conjectured equality fails at $\tau = 1/2$.
\end{proof}

\begin{remark}[Second, independent reason at the same $\tau$]
Writing $v = 1/\tau$, the set $W(\tau)$ is the ``exact order $v$'' set
$\{\alpha : \forall \varepsilon > 0,\ \|q\alpha\| < \varepsilon q^{-v}
\text{ i.m.}\}$. At $\tau = 1/2$ we have $v = 2$, and the Khintchine
divergence test asks for the divergence of
\[
\sum_{q} q \cdot \varepsilon q^{-v} \;=\; \varepsilon \sum_q q^{1-2}
\;=\; \varepsilon \sum_q \frac1q,
\]
the harmonic series, which diverges. Hence $W(1/2)$ has \emph{full
Lebesgue measure}, so its packing dimension equals $1$ — already
different from the conjectured $4/3$ even ignoring the impossibility
argument above.
\end{remark}

\section{Boundary of the attack}

For rational $\tau = p/q > 0$ put $v = q/p$. The conjectured value is
$2/(1+\tau) = 2q/(p+q)$.

\begin{itemize}
\item \textbf{$0 < \tau < 1$ (i.e.\ $q > p$):} the conjectured value
$2q/(p+q)$ exceeds $1$, because $2q > p+q \iff q > p$. The conjecture is
impossible on the entire range $0 < \tau < 1$ by the monotonicity argument;
this includes the primary counterexample $\tau = 1/2$.
\item \textbf{$0 < \tau < 1/2$ (i.e.\ $v > 2$):} here $W(\tau)$ is a null
set and the Jarn\'ik--Besicovitch theory gives
$\dim_H W(\tau) = \dim_P W(\tau) = 2/(v+1) = 2\tau/(1+\tau)$,
which equals the conjectured $2/(1+\tau)$ only at $\tau = 1$. E.g.\ at
$\tau = 1/4$ the true value is $2/5$ while the conjecture claims $8/5$.
\item \textbf{$\tau \ge 1/2$ (i.e.\ $v \le 2$):} Khintchine's theorem gives
full measure, so $\dim_P W(\tau) = 1$, while the conjectured $2/(1+\tau)$
differs from $1$ at every $\tau \neq 1$; e.g.\ at $\tau = 2$ the conjecture
claims $2/3$ but the true value is $1$.
\item \textbf{$\tau = 1$ (the only non-false point):} the conjectured value
is $1$, matching the full-measure value.
\end{itemize}

Thus the conjecture fails for every $\tau > 0$ except the trivial point
$\tau = 1$, and it is outright impossible (value exceeding the ambient
dimension) for all $0 < \tau < 1$. The discrepancy stems from the formula
$2/(1+\tau)$ being the classical Jarn\'ik exponent for the
normalization $q^{\tau}\|q\alpha\|$, whereas the conjecture's
definition uses $q^{1/\tau}$ — the two parametrizations are confused.

\section{Verification}

The arithmetic core of the attack is machine-checked in core Lean~4
(no Mathlib), in the \texttt{lean4/} directory:
\texttt{Tlmc325.attack} proves $\texttt{gtFrac}\ 4\ 3\ 1\ 1 = \texttt{true}$
(i.e.\ $4/3 > 1$), \texttt{Tlmc325.claimed\_exceeds\_ambient} proves the
general $2q/(p+q) > 1$ for $q > p$, and two mismatch theorems record the
true values at $\tau = 2$ and $\tau = 1/4$. Every declaration is proved
with zero axioms (\texttt{Check.lean} prints \emph{does not depend on any
axioms} for all of them) and no \texttt{sorry}. The script
\texttt{reproduce.py} recomputes all numerics standalone with exact
rational arithmetic.

\end{document}
